\documentclass{article}
\usepackage[utf8]{inputenc}
\usepackage{graphicx}
\usepackage{hyperref}
\usepackage{listings}
\usepackage{xcolor}

\title{Computer Networking Mastery}
\author{THE GATEHUB}
\date{\today}

\begin{document}
\maketitle

\learninguniverse{
title={Computer Networking Mastery},
description={A complete course covering networking fundamentals, protocols, security, and hands-on projects.},
difficulty={Intermediate},
estimatedHours={40},
skills={TCP/IP,DNS,HTTP,Routing,Security},
category={Computer Science}
}

\track{
title={Networking Foundations},
description={Core concepts every network engineer must know},

\module{
title={Fundamentals},
description={OSI, TCP/IP, and addressing},

\lesson{title={Lesson 1: Introduction to Networking}}

\overviewmarkdown={
Welcome to Computer Networking Mastery. In this course you will learn how data travels across networks, from physical cables to application protocols.
}

\note{text={Networking connects devices so they can share resources and communicate.}}

\theory{
title={What is a Network?},
body={A computer network is a group of interconnected devices that exchange data using agreed protocols. Networks range from small LANs to the global Internet.}
}

\keypoints{text={Nodes, Links, Protocols, Packets, Bandwidth}}

\quiz{
title={Intro Quiz},
question={What does LAN stand for?},
optionA={Large Area Network},
optionB={Local Area Network},
optionC={Linked Access Node},
optionD={Layered Application Network},
correct={B},
explanation={LAN means Local Area Network — a network in a limited geographic area.}
}

\checkpoint{title={Lesson 1 complete}}

\lesson{title={Lesson 2: The OSI Model}}

\overview{text={This lesson introduces the seven-layer OSI model — the conceptual framework used to understand how network protocols interact.}}

\theory{
title={Seven Layers},
body={The OSI model has 7 layers:
1. Physical
2. Data Link
3. Network
4. Transport
5. Session
6. Presentation
7. Application}
}

\image{
path={https://upload.wikimedia.org/wikipedia/commons/thumb/8/8f/Osi_model.svg/640px-Osi_model.svg.png},
caption={OSI Model diagram},
alt={OSI seven layers}
}

\video{
type={youtube},
url={https://www.youtube.com/watch?v=vvzaS8wca4w},
title={OSI Model Explained}
}

\practice{
title={Identify the Layer},
language={python},
startercode={
# Which OSI layer handles routing?
layer = "Network"
print(layer)
},
expectedoutput={Network}
}

\lesson{title={Lesson 3: TCP/IP Stack}}

\theory{
title={TCP vs UDP},
body={TCP provides reliable, ordered delivery with error checking. UDP is faster but does not guarantee delivery — used for streaming and DNS queries.}
}

\codeexample{
language={python},
code={
import socket
s = socket.socket(socket.AF_INET, socket.SOCK_STREAM)
print("TCP socket created")
},
output={TCP socket created}
}

\quiz{title={TCP/IP Quiz}}

\question{text={Which protocol is connection-oriented?}, explanation={TCP establishes a connection before data transfer.}}
\option{text={TCP}, iscorrect={true}}
\option{text={UDP}, iscorrect={false}}
\option{text={ICMP}, iscorrect={false}}
\option{text={ARP}, iscorrect={false}}

\question{text={What does IP stand for?}}
\option{text={Internet Protocol}, iscorrect={true}}
\option{text={Internal Process}, iscorrect={false}}

}

\module{
title={Protocols and Applications},
description={DNS, HTTP, and application-layer protocols},

\lesson{title={Lesson 4: IP Addressing}}

\tip{text={Remember: IPv4 uses 32 bits; IPv6 uses 128 bits.}}

\practice{
title={Subnet Practice},
language={python},
startercode={
# How many hosts in a /24 network?
hosts = 2**8 - 2
print(hosts)
},
expectedoutput={254},
solution={
hosts = 2**8 - 2
print(hosts)
}
}

\quiz{
question={How many bits in an IPv4 address?},
optionA={16},
optionB={32},
optionC={64},
optionD={128},
correct={B},
explanation={IPv4 addresses are 32 bits, written as four octets.}
}

\discussion{prompt={Why do we still use IPv4 if IPv6 exists?}}

\lesson{title={Lesson 5: DNS and HTTP}}

\theory{
title={DNS Resolution},
body={DNS translates human-readable domain names into IP addresses. A query goes: browser → resolver → root → TLD → authoritative server.}
}

\video{
type={youtube},
url={https://www.youtube.com/watch?v=uOfonONtIuk},
title={How DNS Works}
}

\quiz{
title={Web Protocols},
question={Which port does HTTPS use by default?},
optionA={80},
optionB={443},
optionC={8080},
optionD={22},
correct={B}
}

\resource{
type={website},
title={MDN HTTP Guide},
url={https://developer.mozilla.org/en-US/docs/Web/HTTP}
}

\lesson{title={Lesson 6: Network Security}}

\warning{text={Never expose admin panels to the public internet without authentication.}}

\theory{
title={Common Threats},
body={Malware, phishing, DDoS, man-in-the-middle attacks, and SQL injection are common threats. Defense in depth uses firewalls, encryption, and access controls.}
}

\quiz{
title={Security Quiz},
question={Which are valid security practices? Select all that apply.},
type={multiple},
optionA={Use HTTPS},
optionB={Share passwords in chat},
optionC={Enable 2FA},
optionD={Disable firewalls},
correct={A,C},
explanation={HTTPS and 2FA improve security. Never share passwords or disable firewalls.}
}

\checkpoint{title={Module 2 complete}}

}

\module{
title={Projects and Certification},
description={Capstone project and final assessment},

\lesson{title={Lesson 7: Packet Analyzer Project}}

\project{
title={Build a Packet Logger},
description={Create a simple network packet analyzer.},
instructions={Use Python scapy or socket to capture and display packet headers. Submit your Colab notebook.},
submissionType={colab}
}

\colab{url={https://colab.research.google.com/github/googlecolab/colabtools/blob/master/notebooks/Welcome_To_Colaboratory.ipynb}}

\github{url={https://github.com/example/network-packet-logger}}

\assignment{
title={Written Report},
instructions={Write a 2-page report explaining your packet analyzer design.},
duedate={2026-12-31},
points={100}
}

\lesson{title={Lesson 8: Network Troubleshooting}}

\summary{text={You learned ping, traceroute, netstat, and Wireshark for diagnosing network issues.}}

\practice{
title={Ping Simulation},
language={python},
startercode={
import subprocess
result = "PING 8.8.8.8: OK"
print(result)
},
expectedoutput={PING 8.8.8.8: OK}
}

\download{
title={Troubleshooting Cheatsheet},
url={https://example.com/network-cheatsheet.pdf}
}

\lesson{title={Lesson 9: Final Review}}

\certificatecriteria{
Complete all lessons, score 70% or higher on quizzes, and submit the capstone project to earn your certificate.
}

\finalexam{
title={Computer Networking Final Exam},
duration={60 minutes},
description={Comprehensive assessment covering OSI, TCP/IP, DNS, HTTP, and security topics.}
}

\quiz{
title={Final Exam Sample},
question={Which layer of the OSI model does a router primarily operate at?},
optionA={Physical},
optionB={Data Link},
optionC={Network},
optionD={Application},
correct={C},
explanation={Routers operate at Layer 3 (Network), making routing decisions based on IP addresses.}
}

}

}

\end{document}
